% ===========================================================================
\section{Properties of Strategies}
\label{sec:properties}

The archetypes of Section~\ref{sec:archetypes} differ along a small
number of recurring dimensions.
This section makes those dimensions precise.
Throughout, fix a pool $\mathcal{S}$ of $n$ strategies and write
$d(s, r) \in \{\textup{\textsc{share}}, \textup{\textsc{take}}\}$ for
the oracle's decision when $s$ dictates to $r$; the decision profile
is the restriction of $d$ to $\mathcal{S}^2$.
Properties defined below are relative to this profile.
(When the oracle is sampled, $d(s,r)$ may be read as the majority
decision, or the definitions restated in expectation.)

\subsection{Score decomposition and the exchange rate}

Two constants govern every trade-off in the game:
the dictator's cost of sharing and the recipient's gain from being
shared with,
\begin{equation}
  \delta = u(E) - u(E/2), \qquad g = u(E/2).
  \label{eq:delta-g}
\end{equation}
Let $k_{\mathrm{give}}(s) = |\{r : d(s,r) = \textup{\textsc{share}}\}|$
and
$k_{\mathrm{recv}}(s) = |\{r : d(r,s) = \textup{\textsc{share}}\}|$.
A strategy's tournament score decomposes as
\begin{equation}
  \mathrm{Score}(s)
  = n\, u(E) \;-\; \delta\, k_{\mathrm{give}}(s)
  \;+\; g\, k_{\mathrm{recv}}(s),
  \label{eq:score-decomposition}
\end{equation}
so any comparison between strategies reduces to
$g\,\Delta k_{\mathrm{recv}} \gtrless \delta\,\Delta k_{\mathrm{give}}$.
The \emph{exchange rate}
\begin{equation}
  \rho = \frac{g}{\delta}
       = \frac{\ln(1 + E/2)}{\ln(1+E) - \ln(1 + E/2)}
  \label{eq:exchange-rate}
\end{equation}
measures how many shares given are offset by one share received;
for $E = 60$, $\rho \approx 5.07$.
Conditional cooperation is cheap in this game: a strategy can afford
to share with up to $\lfloor \rho \rfloor$ opponents for each
additional dictator it induces to share with it.
The exchange rate recurs throughout: it sets the altruism threshold
$\lambda^*= 1/\rho$ of Section~\ref{sec:bayesian}, and the
break-even point for exploiting unconditional cooperators
(Equation~\eqref{eq:protection-breakeven} below).

\subsection{First-order properties}

\begin{definition}[Unconditionality]
Strategy $s$ is \emph{unconditional} if $d(s, r)$ does not depend on
$r$.  The two constant strategies are the unconditional cooperator
(\textbf{Generous}) and the unconditional defector (\textbf{Selfish});
all other strategies are \emph{conditional}.
\end{definition}

\begin{definition}[Self-cooperation]
Strategy $s$ is \emph{self-cooperating} if
$d(s, s) = \textup{\textsc{share}}$.
\end{definition}

Since the round robin includes self-play, failing self-cooperation
costs $2u(E/2) - u(E) = g - \delta > 0$ outright.
Self-play is also the degenerate case of decision-process
correlation: a strategy's copy decides identically, so
self-cooperation is the minimal form of the FDT-style reasoning
discussed in Section~\ref{sec:related}.

\begin{definition}[Reciprocity]
Strategy $s$ is \emph{reciprocal} if $d(s, r) = d(r, s)$ for all
$r \in \mathcal{S}$: it treats every opponent exactly as that
opponent treats it (\textbf{Mirror}).
\end{definition}

\begin{definition}[Behavioral equivalence and extensionality]
Strategies $r, r'$ are \emph{behaviorally equivalent} if they have
identical rows and columns in the decision profile:
$d(r, x) = d(r', x)$ and $d(x, r) = d(x, r')$ for all $x$.
Strategy $s$ is \emph{extensional} if $d(s, r) = d(s, r')$ whenever
$r$ and $r'$ are behaviorally equivalent.
\end{definition}

An intensional strategy reacts to features of the opponent's artifact
that have no behavioral consequence --- self-labels, style, or
keyword handshakes.
The no-collusion rule of our tournaments is in effect a ban on one
intensional class (secret handshakes).

\begin{definition}[Exploitability]
Strategy $s$ is \emph{exploited by} $r$ if
$d(s, r) = \textup{\textsc{share}}$ and
$d(r, s) = \textup{\textsc{take}}$: $s$ subsidizes an opponent that
takes from it.
Relative to the profile, call $r$ an \emph{exploiter} if it exploits
some strategy in $\mathcal{S}$, and call a strategy that shares with
an exploiter a \emph{subsidizer}.
\end{definition}

\begin{definition}[First-order deterrence]
Strategy $s$ is a \emph{first-order deterrent} if
$d(s, r) = \textup{\textsc{take}}$ for every exploiter $r$.
\end{definition}

Every conditional archetype of Section~\ref{sec:archetypes} is a
first-order deterrent with respect to the unconditional defector.
By Equation~\eqref{eq:score-decomposition}, first-order deterrence is
what makes defection unprofitable: a defector's only revenue beyond
the baseline $n\,u(E)$ is $g$ per subsidizer, so in a pool of
first-order deterrents it earns the pool minimum.

\subsection{Second-order norms: deterrence and protection}
\label{sec:second-order}

First-order deterrence leaves a gap.
In a pool of conditional cooperators with defectors absent, an
unconditional cooperator earns exactly what the conditional
cooperators earn --- every entry of its row and column is
\textup{\textsc{share}} --- so nothing maintains conditionality, and
each unconditional cooperator raises the revenue of a would-be
defector by $g$.
Two archetypes in our pool respond to this second-order problem, in
opposite directions.

\begin{definition}[Second-order deterrence]
Strategy $s$ is a \emph{second-order deterrent} if it is a
first-order deterrent and additionally
$d(s, r) = \textup{\textsc{take}}$ for every subsidizer $r$.
(\textbf{Cooperation coalition}.)
\end{definition}

\begin{proposition}
\label{prop:drift}
Let $\mathcal{S}$ contain a second-order deterrent $s_2$ and a set
$C$ of mutually sharing conditional cooperators that share with
$s_2$ and with the unconditional cooperator $\bar{c}$, with $s_2$
sharing with $C$.
Then $\bar{c}$ earns strictly less than every member of $C$, by at
least $g$.
\end{proposition}

\begin{proof}
$\bar{c}$ and any $c \in C$ have $k_{\mathrm{give}}$ differing by at
most the defectors and subsidizers that $\bar{c}$ shares with and $c$
does not, each such round \emph{raising}
$\mathrm{Score}(c) - \mathrm{Score}(\bar{c})$ by $\delta$ via
Equation~\eqref{eq:score-decomposition}.
On the receiving side, $s_2$ takes from $\bar{c}$ (a subsidizer, as
$\bar{c}$ shares with every exploiter) and shares with $c$, so
$k_{\mathrm{recv}}(\bar{c}) \le k_{\mathrm{recv}}(c) - 1$.
Hence
$\mathrm{Score}(c) - \mathrm{Score}(\bar{c}) \ge g$.
\end{proof}

Second-order deterrence thus closes the neutral-drift channel: it
makes unconditional cooperation strictly suboptimal even when
defectors are extinct, stabilizing the all-conditional configuration.
The cost is that the punishment falls on the most cooperative
strategies in the pool whenever they do appear.

\begin{definition}[Protection]
Strategy $s$ is \emph{protective} if
$d(s, r) = d(r, \bar{c})$ for all $r$, where $\bar{c}$ is the
unconditional cooperator: $s$ treats each opponent as that opponent
would treat an innocent.
(\textbf{Chivalry}.)
\end{definition}

A protective strategy addresses the same second-order problem by
making the \emph{exploitation} of unconditional cooperators costly
rather than their existence.
Exploiting one unconditional cooperator harvests $\delta$; each
protector responds by withholding its share, costing $g$.
Harvesting is therefore profitable only if the pool contains enough
innocents per protector:
\begin{equation}
  \#\{\text{unconditional cooperators harvested}\}
  \;>\; \rho \cdot \#\{\text{protectors}\}.
  \label{eq:protection-breakeven}
\end{equation}

The two norms are mutually antagonistic.
An unconditional cooperator is a subsidizer, so a second-order
deterrent takes from it; a protector, applying the innocent test to
the second-order deterrent, finds
$d(s_2, \bar{c}) = \textup{\textsc{take}}$ and takes from $s_2$ in
turn.
Each second-order norm punishes precisely the behavior the other
mandates toward the unconditionally generous.
Both stabilize full cooperation among conditional strategies; they
differ in whether unconditional cooperators end up extinct or
protected.
This antagonism is not a curiosity: it generates the two families of
equilibria --- enforcement and protection --- that appear in the
Bayesian population analysis of Section~\ref{sec:bayesian}, and
Equation~\eqref{eq:protection-breakeven} determines which family a
given pool composition favors.

\subsection{Legibility and well-foundedness:
open strategies versus open source}
\label{sec:properties-oracle}

Two further properties matter greatly in practice but are
deliberately excluded from the list above: they are properties of the
strategy--oracle pair, not of the decision profile, and they are
where the LLM-adjudicated tournament differs most from its
predecessors in which formal programs analyze other programs
\citep{tennenholtz2004program, ospd2013lesswrong}.

\paragraph{Legibility.}
By Equation~\eqref{eq:score-decomposition}, a strategy's score
depends on $k_{\mathrm{recv}}$ --- on how \emph{other} dictators
classify its artifact --- and the exchange rate makes errors on this
side roughly $\rho$ times as costly as errors the strategy makes as
dictator.
A strategy is \emph{legible} to a class of conditional dictators if
they classify it as the strategy's behavior warrants.
In program-equilibrium settings legibility is provability:
cooperation obtains only when a machine-checkable proof about the
opponent's source exists, so behaviorally equivalent but
syntactically distinct programs defeat recognition, and cooperation
is brittle by construction.
LLM adjudication replaces proof with judgment: legibility becomes
graded and robust to paraphrase, at the price of classification
noise.
Misreadings replace impossibility results --- an aggressive
classification rule that would simply fail to find proofs in the
formal setting instead occasionally misfires against cooperative
opponents in ours, a failure mode we return to in
Section~\ref{sec:discussion}.

\paragraph{Well-foundedness.}
Mutual conditionality is self-referential.
When two reciprocal strategies meet, ``share iff they would share
with me'' admits two consistent resolutions --- mutual
\textup{\textsc{share}} and mutual \textup{\textsc{take}} --- and no
finite unwinding selects between them.
A strategy is \emph{well-founded} if its decision against every
opponent in the pool is determined by a finite regress.
In the formal-program setting this is the L\"obian obstacle, and
resolving it requires provability-logic constructions engineered for
the purpose \citep{barasz2014robust} or simulation with stochastic
grounding \citep{oesterheld2019robust}.
In our setting the oracle itself acts as the fixed-point selector,
and empirically selects the cooperative resolution --- a disposition
of the arena, not a property of any strategy.
Strategies can also restore well-foundedness on their own:
\textbf{Universalizability} grounds the regress at depth one by
evaluating a population of the opponent's copies, and
\textbf{Chivalry} grounds it by substituting the unconditional
cooperator into the recursion, which is what makes the protective
norm of Section~\ref{sec:second-order} implementable at all.
